% ============================================================
% 纯答案册·课时16-2.7.2抛物线性质 —— 工具/答案抽册器.py 抽册生成件（勿手改；第三档＝纯答案单独成册）
% 源件（只读）：C:/提示词/工作区/M3-第2章量产0913/成卷/导学件/课时16-2.7.2抛物线性质/main.tex（md5 c0c21bfcea271ce455db3d166fb3f326）
%% 口径：题面不重印；逐题＝题号(印面号同正册)＋[答案]＋[详解]，值/详解照源件逐字
%       （钉值门硬断言：册值≡正册件内值≡值台账值tex，逐字全等）。册序＝正册装配序＝印面号 1..50 升序（同序同号）；键序断言基准＝题面库冻结 manifest canonical 键序。
%% 三档导出：整体 main-true／纯题 main-false（在产）｜本册＝纯答案册（本件）
%% 双档：ansbook-true.tex＝详解印本；ansbook-false.tex＝纯值速查本（详解不印）
% 编译：xelatex <壳> ×2，cwd＝本目录；sty＝qp-m3.sty 本地挂载（md5 7c3930362be8a0a2bdf21bbf8ac16573，锁校验）
% ============================================================
\documentclass[fontset=none]{ctexart}
\usepackage{qp-m3}
% —— 印面逐字符号路由补钉（承源件；− 走 TNR、▱ NSC 子块；禁改 sty 保 md5） ——
\xeCJKDeclareCharClass{Default}{"2212}%
\setCJKfamilyfont{nscblk}[Path=C:/提示词/工作区/字替对照-0909/variantF/fonts/,BoldFont=NSC-w600.ttf]{NSC-w500.ttf}%
\xeCJKDeclareSubCJKBlock{nscblk}{"25B1}%
\newcommand{\pxparallelogram}{{\CJKfamily{nscblk}▱}}%
% —— 断行微调（承母版实测档，三零门承重；\emergencystretch 不另设） ——
\xeCJKsetup{CJKglue={\hskip 0.10em plus 0.08em minus 0.05em}}%
% —— 纯答案册全线括线（长详解可跨栏断，承练习件全线括线口径；灰底盒不可跨栏断） ——
\ansblockgrayfalse
% —— 双档开关（false 档＝纯值速查：答案印、详解不印；件面开关，不动 sty 本体） ——
\newif\ifansbookdetail
\ifdefined\ansbookdetail
  \ifnum\ansbookdetail=0 \ansbookdetailfalse\else\ansbookdetailtrue\fi
\else
  \ansbookdetailtrue
\fi
% —— 页脚件名（承源件章名；件名＝<件型>答案册） ——
\renewcommand{\qpzhangming}{第二章\quad 平面解析几何}%
\renewcommand{\qpjianming}{导学件答案册}%

\begin{document}
\zhangtitle{第二章\quad 平面解析几何}
\jietitle{2.7.2 抛物线的几何性质}
\par\glueguard{1}\addvspace{2pt}\noindent\biaoqian{【纯答案册】}{\fangsong 题面不重印\quad 印面号与正册同号同序}\par\addvspace{2pt}

\begin{multicols}{2}
\emergencystretch=1em

\begin{ansblock}[2章-练-课时16-04]
% ans:2章-练-课时16-04
\ansitem{1}{1}
\ifansbookdetail
\ansnote{详解}{\(2p=8\)、\(p=4\)，焦点 \(F(2,0)\)．由点到直线的距离公式，\(d=|2-\sqrt{3}\times 0|/\sqrt{1^{2}+(\sqrt{3})^{2}}=2/\sqrt{4}=1\)．}
\fi
\end{ansblock}

\begin{ansblock}[2章-练-课时16-03]
% ans:2章-练-课时16-03
\ansitem{2}{(1) 焦点(3/2,0)，准线 x=−3/2；(2) 焦点(0,−5/2)，准线 y=5/2；(3) 焦点(0,2/3)，准线 y=−2/3；(4) 焦点(−a/4,0)，准线 x=a/4}
\ifansbookdetail
\ansnote{详解}{(1) \(y^{2}=6x\)：\(2p=6\)、\(p=3\)，开口向右，焦点 \((p/2,0)=(3/2,0)\)，准线 \(x=-3/2\)．(2) \(x^{2}=-10y\)：\(2p=10\)、\(p=5\)，开口向下，焦点 \((0,-5/2)\)，准线 \(y=5/2\)．(3) \(x^{2}=(8/3)y\)：\(2p=8/3\)、\(p=4/3\)，开口向上，焦点 \((0,2/3)\)，准线 \(y=-2/3\)．(4) \(y^{2}=-ax\)（\(a\neq 0\)）：一次项系数为 \(-a\)，无论 \(a>0\)（开口向左）或 \(a<0\)（开口向右），焦点均为 \((-a/4,0)\)、准线均为 \(x=a/4\)（\(a>0\) 时焦点在 \(x\) 轴负半轴，\(a<0\) 时在正半轴，公式统一）．}
\fi
\end{ansblock}

\begin{ansblock}[2章-练-课时16-01]
% ans:2章-练-课时16-01
\ansitem{3}{AD}
\ifansbookdetail
\ansnote{详解}{对照标准形 \(y^{2}=-2px\)（\(p>0\)）：\(2p=2\)、\(p=1\)，一次项系数为负，开口向左，A 正确；焦点为 \((-p/2,0)=(-1/2,0)\)，B 所给 \((-1,0)\) 错（把 \(p\) 误作 \(p/2\)）；准线为 \(x=p/2=1/2\)，C 所给 \(x=1\) 错；对称轴为 \(x\) 轴，D 正确．故选：AD．}
\fi
\end{ansblock}

\begin{ansblock}[2章-练-课时16-05]
% ans:2章-练-课时16-05
\ansitem{4}{6}
\ifansbookdetail
\ansnote{详解}{\(y^{2}=8x\) 中 \(2p=8\)、\(p=4\)，准线 \(x=-2\)．点 \(P\) 到 \(y\) 轴距离为 4，而抛物线上点横坐标 \(x\geqslant 0\)，故 \(x_{P}=4\)．由定义，\(P\) 到焦点的距离等于 \(P\) 到准线的距离：\(|PF|=x_{P}+p/2=4+2=6\)．（校验：\(P(4,\pm 4\sqrt{2})\)，\(|PF|=\sqrt{(4-2)^{2}+32}=\sqrt{36}=6\) \ding{51}．）}
\fi
\end{ansblock}

\begin{ansblock}[2章-练-课时16-02]
% ans:2章-练-课时16-02
\ansitem{5}{17/16}
\ifansbookdetail
\ansnote{详解}{化标准形：\(y=4x^{2}\) 即 \(x^{2}=(1/4)y\)，故 \(2p=1/4\)、\(p=1/8\)，焦点 \(F(0,1/16)\)，准线 \(y=-1/16\)．由抛物线定义，点 \(M\) 到焦点的距离等于它到准线的距离，\(|MF|=y_{M}+p/2=1+1/16=17/16\)．（回代校验：\(x_{M}^{2}=1/4\)、\(M(\pm 1/2,1)\)，\(|MF|^{2}=x_{M}^{2}+(y_{M}-1/16)^{2}=1/4+225/256=289/256\)，\(|MF|=17/16\) \ding{51}．）}
\fi
\end{ansblock}

\begin{ansblock}[2章-练-课时16-06]
% ans:2章-练-课时16-06
\ansitem{6}{M(3p/2, ±√3 p)}
\ifansbookdetail
\ansnote{详解}{设 \(M(x_{0},y_{0})\)，\(x_{0}\geqslant 0\)．由定义，\(|MF|\) 等于 \(M\) 到准线 \(x=-p/2\) 的距离，即 \(x_{0}+p/2=2p\)，得 \(x_{0}=3p/2\)．代入方程：\(y_{0}^{2}=2p\cdot(3p/2)=3p^{2}\)，\(y_{0}=\pm\sqrt{3}p\)．故 \(M(3p/2,\sqrt{3}p)\) 或 \(M(3p/2,-\sqrt{3}p)\)．}
\fi
\end{ansblock}

\begin{ansblock}[2章-练-课时16-07]
% ans:2章-练-课时16-07
\ansitem{7}{y²=−4x 或 x²=√2 y}
\ifansbookdetail
\ansnote{详解}{\(P(-2,2\sqrt{2})\) 在第二象限，开口只可能向左或向上．情形一，对称轴为 \(x\) 轴：设 \(y^{2}=-2px\)（\(p>0\)），代入得 \((2\sqrt{2})^{2}=-2p\cdot(-2)=4p\)、\(p=2\)，即 \(y^{2}=-4x\)．情形二，对称轴为 \(y\) 轴：设 \(x^{2}=2py\)（\(p>0\)），代入得 \(4=2p\cdot 2\sqrt{2}\)、\(2p=\sqrt{2}\)，即 \(x^{2}=\sqrt{2}y\)．（回验：\(8=-4\times(-2)\) \ding{51}；\(4=\sqrt{2}\times 2\sqrt{2}\) \ding{51}．）}
\fi
\end{ansblock}

\begin{ansblock}[2章-练-课时16-08]
% ans:2章-练-课时16-08
\ansitem{8}{p=4}
\ifansbookdetail
\ansnote{详解}{椭圆中 \(c^{2}=6-2=4\)，右焦点 \((2,0)\)．抛物线焦点 \((p/2,0)\)，由 \(p/2=2\) 得 \(p=4\)．}
\fi
\end{ansblock}

\begin{ansblock}[2章-练-课时16-10]
% ans:2章-练-课时16-10
\ansitem{9}{C}
\ifansbookdetail
\ansnote{详解}{顶点在原点、对称轴为 \(x\) 轴，方程形如 \(y^{2}=\pm 2px\)（\(p>0\)）．通径（过焦点且垂直于对称轴的弦）长为 \(2p\)，由 \(2p=8\) 得 \(p=4\)，故方程为 \(y^{2}=8x\)（焦点在 \(x\) 轴正半轴）或 \(y^{2}=-8x\)（焦点在负半轴），两向各一，选 C．（排除 D：其对称轴为 \(y\) 轴．）}
\fi
\end{ansblock}

\begin{ansblock}[2章-练-课时16-09]
% ans:2章-练-课时16-09
\ansitem{10}{A（1/3）}
\ifansbookdetail
\ansnote{详解}{\(y^{2}=x\) 中 \(2p=1\)，焦点 \(F(1/4,0)\)，准线 \(x=-1/4\)．直线 \(l\)：\(y=\sqrt{3}(x-1/4)\)．代入 \(y^{2}=x\) 得 \(3(x-1/4)^{2}=x\)，整理为 \(48x^{2}-40x+3=0\)，故 \(x_{1}+x_{2}=40/48=5/6\)，\(x_{1}x_{2}=3/48=1/16\)．由定义，\(|AF|=x_{1}+1/4\)、\(|BF|=x_{2}+1/4\)，\(|AF|\cdot|BF|=x_{1}x_{2}+(1/4)(x_{1}+x_{2})+1/16=1/16+5/24+1/16=16/48=1/3\)．故选 A．}
\fi
\end{ansblock}

\begin{ansblock}[2章-练-课时16-11]
% ans:2章-练-课时16-11
\ansitem{11}{B}
\ifansbookdetail
\ansnote{详解}{\(y^{2}=4x\) 的焦点 \(F(1,0)\)，准线 \(x=-1\)，即直线 \(x+1=0\)．分别过 \(A\)、\(B\)、\(M\) 作准线的垂线，垂足为 \(C\)、\(D\)、\(N\)．由定义，\(|AC|=|AF|\)、\(|BD|=|BF|\)，故 \(|AC|+|BD|=|AF|+|BF|=|AB|=8\)．又 \(AC\parallel MN\parallel BD\) 且 \(M\) 为 \(AB\) 中点，\(MN\) 为直角梯形 \(ABDC\) 的中位线，\(|MN|=(|AC|+|BD|)/2=4\)，即点 \(M\) 到准线 \(x=-1\) 的距离为 4．故选：B．}
\fi
\end{ansblock}

\begin{ansblock}[2章-练-课时16-12]
% ans:2章-练-课时16-12
\ansitem{12}{P(−2,1/2)（最小值 6）}
\ifansbookdetail
\ansnote{详解}{\(x^{2}=8y\) 的准线为 \(l\)：\(y=-2\)．设 \(P\) 在 \(l\) 上的射影为 \(Q\)，过 \(A\) 作 \(AB\perp l\) 于 \(B\)，则 \(|AB|=y_{A}-(-2)=4+2=6\)．由定义 \(|PF|=|PQ|\)，故 \(|PF|+|PA|=|PQ|+|PA|\geqslant|AQ|\geqslant|AB|=6\)：第一个不等号当 \(A\)、\(P\)、\(Q\) 共线（\(P\) 与 \(A\) 同取横坐标 \(-2\)）时取等，第二个不等号当 \(Q\) 与 \(B\) 重合时取等，两条件同时满足．把 \(x=-2\) 代入 \(x^{2}=8y\) 得 \(y=1/2\)，故所求点 \(P(-2,1/2)\)，此时 \((|PF|+|PA|)_{\min}=|AB|=y_{A}+2=6\)．}
\fi
\end{ansblock}

\begin{ansblock}[2章-练-课时16-13]
% ans:2章-练-课时16-13
\ansitem{13}{|MA| 的最小值为 4，此时 M(0,0)}
\ifansbookdetail
\ansnote{详解}{设 \(M(x,y)\)，则 \(x\geqslant 0\)，\(y^{2}=16x\)．\par\noindent \(|MA|^{2}=(x-4)^{2}+y^{2}=(x-4)^{2}+16x=x^{2}+8x+16=(x+4)^{2}\)．\par\noindent 由 \(x\geqslant 0\) 得 \((x+4)^{2}\geqslant 16\)，当且仅当 \(x=0\) 时等号成立，此时 \(y=0\)，\(|MA|_{\min}=4\)，\(M(0,0)\)．（注：\(2p=16\)、\(p/2=4\)，\(A(4,0)\) 恰为该抛物线的焦点，\(|MA|=|MF|=x+4\geqslant 4\) 与配方路径互证．）}
\fi
\end{ansblock}

\begin{ansblock}[2章-练-课时16-14]
% ans:2章-练-课时16-14
\ansitem{14}{4√3 p}
\ifansbookdetail
\ansnote{详解}{设正三角形 \(OAB\) 的顶点 \(A(x_{1},y_{1})\)、\(B(x_{2},y_{2})\) 在抛物线上，则 \(y_{1}^{2}=2px_{1}\)、\(y_{2}^{2}=2px_{2}\)．由 \(|OA|=|OB|\) 得 \(x_{1}^{2}+y_{1}^{2}=x_{2}^{2}+y_{2}^{2}\)，即 \((x_{1}^{2}-x_{2}^{2})+2p(x_{1}-x_{2})=0\)，\((x_{1}-x_{2})(x_{1}+x_{2}+2p)=0\)．因 \(x_{1}\geqslant 0\)、\(x_{2}\geqslant 0\) 且 \(2p>0\)（若 \(x_{1}\)、\(x_{2}\) 之一为 0 则三角形退化），括号第二因子为正，故 \(x_{1}=x_{2}\)，从而 \(y_{2}=-y_{1}\neq 0\)，线段 \(AB\) 关于 \(x\) 轴对称．\(x\) 轴垂直平分 \(AB\) 且平分 \(\angle AOB\)，\(\angle AOx=30^{\circ}\)，故 \(y_{1}/x_{1}=\tan 30^{\circ}=\sqrt{3}/3\)、\(x_{1}=\sqrt{3}y_{1}\)．代入 \(y_{1}^{2}=2p\cdot\sqrt{3}y_{1}\) 得 \(y_{1}=2\sqrt{3}p\)，边长 \(|AB|=2y_{1}=4\sqrt{3}p\)．}
\fi
\end{ansblock}

\begin{ansblock}[2章-练-课时16-15]
% ans:2章-练-课时16-15
\ansitem{15}{(1) |AB|=4p；(2) 定值 cos∠AOB=−3√41/41（与 p 无关）}
\ifansbookdetail
\ansnote{详解}{(1) 过 \(F(p/2,0)\)、倾斜角 \(\pi/4\) 的直线为 \(y=x-p/2\)．代入 \(y^{2}=2px\)：\((x-p/2)^{2}=2px\)，即 \(x^{2}-3px+p^{2}/4=0\)（\(\Delta=9p^{2}-p^{2}=8p^{2}>0\)）．设 \(A(x_{A},y_{A})\)、\(B(x_{B},y_{B})\)，则 \(x_{A}+x_{B}=3p\)，\(x_{A}x_{B}=p^{2}/4\)．由定义（焦半径），\par\noindent \(|AB|=|AF|+|BF|=(x_{A}+p/2)+(x_{B}+p/2)=x_{A}+x_{B}+p=4p\)．\par\noindent (2) 该弦所在直线与 \(x\) 轴交于 \(F\)，两端点 \(A\)、\(B\) 分居 \(x\) 轴两侧，纵坐标一正一负，其积为负：\par\noindent \(y_{A}y_{B}=(x_{A}-p/2)(x_{B}-p/2)=x_{A}x_{B}-(p/2)(x_{A}+x_{B})+p^{2}/4=p^{2}/4-3p^{2}/2+p^{2}/4=-p^{2}\)（与焦点弦结论 \(y_{1}y_{2}=-p^{2}\) 相符，此即符号判据）．\par\noindent 于是 \(x_{A}x_{B}+y_{A}y_{B}=p^{2}/4-p^{2}=-3p^{2}/4\)．又 \(|OA|^{2}=x_{A}^{2}+y_{A}^{2}=x_{A}^{2}+2px_{A}=x_{A}(x_{A}+2p)\)，同理 \(|OB|^{2}=x_{B}(x_{B}+2p)\)．\par\noindent 故 \(|OA|\cdot|OB|=\sqrt{x_{A}x_{B}}\cdot\sqrt{x_{A}x_{B}+2p(x_{A}+x_{B})+4p^{2}}=(p/2)\cdot\sqrt{p^{2}/4+6p^{2}+4p^{2}}=(p/2)\cdot(\sqrt{41}/2)p=\sqrt{41}p^{2}/4\)．\par\noindent 在 \(\triangle AOB\) 中由余弦定理，\(\cos\angle AOB=(|OA|^{2}+|OB|^{2}-|AB|^{2})/(2|OA|\cdot|OB|)\)．\par\noindent 而 \(|OA|^{2}+|OB|^{2}-|AB|^{2}=(x_{A}-x_{B})^{2}+(y_{A}-y_{B})^{2}=2(x_{A}x_{B}+y_{A}y_{B})=-3p^{2}/2\)，\par\noindent 故 \(\cos\angle AOB=(-3p^{2}/4)/(\sqrt{41}p^{2}/4)=-3/\sqrt{41}\)，分母有理化得 \(-3\sqrt{41}/41\)，\(p^{2}\) 全部约去，与 \(p\) 无关．}
\fi
\end{ansblock}

\begin{ansblock}[2章-练-课时16-16]
% ans:2章-练-课时16-16
\ansitem{16}{24}
\ifansbookdetail
\ansnote{详解}{\(x^{2}=4y\) 的焦点 \(F(0,1)\)，准线 \(y=-1\)，抛物线上点之焦半径 \(|AF|=y_{A}+1\)（到准线距离，焦半径公式见例 36）．\(l_{1}\)：\(y=k_{1}x+1\) 代入 \(x^{2}=4y\) 得 \(x^{2}-4k_{1}x-4=0\)（\(\Delta=16k_{1}^{2}+16>0\) 恒成立）．设 \(A(x_{1},y_{1})\)、\(B(x_{2},y_{2})\)，则 \(x_{1}+x_{2}=4k_{1}\)，\(y_{1}+y_{2}=k_{1}(x_{1}+x_{2})+2=4k_{1}^{2}+2\)．于是 \par\noindent \(|AB|=|AF|+|BF|=(y_{1}+1)+(y_{2}+1)=y_{1}+y_{2}+2=4k_{1}^{2}+4\)．\par\noindent 同理 \(|DE|=4k_{2}^{2}+4\)．故 \par\noindent \(|AB|+|DE|=4(k_{1}^{2}+k_{2}^{2})+8\geqslant 8|k_{1}k_{2}|+8=8\times 2+8=24\)，\par\noindent 当且仅当 \(|k_{1}|=|k_{2}|\) 且 \(k_{1}k_{2}=-2\)，即 \((k_{1},k_{2})=(\sqrt{2},-\sqrt{2})\) 或 \((-\sqrt{2},\sqrt{2})\) 时等号成立．所求最小值为 24．}
\fi
\end{ansblock}

\begin{ansblock}[2章-拓-课时16-01]
% ans:2章-拓-课时16-01
\ansitem{17}{B}
\ifansbookdetail
\ansnote{详解}{点在 \(y^{2}=4x\) 上，故 \(x\geqslant 0\) 且 \(y^{2}=4x\)．代入得 \(z=x^{2}+(1/2)\cdot 4x+3=x^{2}+2x+3=(x+1)^{2}+2\)．由 \(x\geqslant 0\)，\(x=0\) 时 \(z\) 最小，最小值为 3．故选：B．}
\fi
\end{ansblock}

\begin{ansblock}[2章-拓-课时16-03]
% ans:2章-拓-课时16-03
\ansitem{18}{C}
\ifansbookdetail
\ansnote{详解}{\(x^{2}=4y\) 的焦点 \(F(0,1)\)，准线 \(y=-1\)．延长 \(QP\) 交准线于 \(M\)，则 \(|PM|=y_{P}+1=|PQ|+1\)；由定义 \(|PM|=|PF|\)，故 \(|PQ|=|PF|-1\)．于是 \par\noindent \(|PA|+|PQ|=|PA|+|PF|-1\geqslant|AF|-1=\sqrt{8^{2}+(7-1)^{2}}-1=10-1=9\)，\par\noindent 当且仅当 \(A\)、\(P\)、\(F\) 三点共线时取等（\(A(8,7)\) 在抛物线外：\(64>4\times 7=28\)，线段 \(AF\) 与抛物线相交，取等可行）．最小值为 9，选 C．}
\fi
\end{ansblock}

\begin{ansblock}[2章-拓-课时16-04]
% ans:2章-拓-课时16-04
\ansitem{19}{D}
\ifansbookdetail
\ansnote{详解}{化标准形 \(x^{2}=-4y\)：焦点 \(F(0,-1)\)，准线 \(l\)：\(y=1\)．第一根式是 \(P\) 到 \(F(0,-1)\) 的距离 \(|PF|\)，第二根式是 \(P\) 到 \(A(4,-5)\) 的距离 \(|PA|\)．由定义，\(|PF|\) 等于 \(P\) 到准线 \(y=1\) 的距离（\(n\leqslant 0\)，即 \(1-n\)），故 \(|PF|+|PA|\) 不小于 \(A\) 到准线的垂线段长 \(1-(-5)=6\)（\(A\) 在抛物线内部：\(4^{2}=16<-4\times(-5)=20\) \ding{51}，\(P\) 取垂线段与抛物线交点即可取等）．最小值 6，选 D．}
\fi
\end{ansblock}

\begin{ansblock}[2章-拓-课时16-02]
% ans:2章-拓-课时16-02
\ansitem{20}{42 或 22}
\ifansbookdetail
\ansnote{详解}{按 \(M\) 与抛物线的位置分类．(1) \(M\) 在抛物线内部或上：\(40^{2}\leqslant 2p\cdot 20\)，即 \(p\geqslant 40\)．设 \(P\) 在准线 \(x=-p/2\) 上的射影为 \(D\)，\(|PF|=|PD|\)，\(M\)、\(P\)、\(D\) 共线时取最小，最小值为 \(M\) 到准线的距离 \(20+p/2=41\)，得 \(p=42\)（\(40\leqslant 42\) \ding{51}）．(2) \(M\) 在抛物线外部：\(p<40\)，\(P\)、\(M\)、\(F\) 共线时取最小，最小值为 \par\noindent \(|MF|=\sqrt{40^{2}+(20-p/2)^{2}}=41\)，\((20-p/2)^{2}=81\)、\(p/2=11\) 或 \(29\)，\(p=22\) 或 \(p=58\)（\(58>40\) 与外部前提矛盾，舍）．\par\noindent 综上 \(p=42\) 或 \(22\)．}
\fi
\end{ansblock}

\begin{ansblock}[2章-拓-课时16-21]
% ans:2章-拓-课时16-21
\ansitem{21}{7/2；P(2,2)}
\ifansbookdetail
\ansnote{详解}{\(A(3,2)\)：\(2^{2}=4<2\times 3=6\)，\(A\) 在抛物线内部．准线 \(l\)：\(x=-1/2\)．设 \(P\) 在 \(l\) 上的射影为 \(Q\)，由定义 \(|PF|=|PQ|\)，\(|PA|+|PF|=|PA|+|PQ|\geqslant A\) 到准线的水平距离 \(=3-(-1/2)=7/2\)，当 \(A\)、\(P\)、\(Q\) 共线（\(y_{P}=y_{A}=2\)）时取等．代入 \(y^{2}=2x\) 得 \(x=2\)，即 \(P(2,2)\)，最小值 \(7/2\)．}
\fi
\end{ansblock}

\begin{ansblock}[2章-拓-课时16-09]
% ans:2章-拓-课时16-09
\ansitem{22}{(0,0)（最小值 8）}
\ifansbookdetail
\ansnote{详解}{设 \(P(x_{0},y_{0})\)，\(y_{0}^{2}=-4x_{0}\)（\(x_{0}\leqslant 0\)）．\par\noindent \(\overrightarrow{AP}\cdot\overrightarrow{BP}=(x_{0}-2)(x_{0}-4)+y_{0}^{2}=x_{0}^{2}-6x_{0}+8-4x_{0}=x_{0}^{2}-10x_{0}+8=(x_{0}-5)^{2}-17\)．\par\noindent 由 \(x_{0}\leqslant 0\)，\((x_{0}-5)^{2}\geqslant 25\) 且随 \(x_{0}\) 增大而减小，故 \(x_{0}=0\) 时 \(\overrightarrow{AP}\cdot\overrightarrow{BP}\) 最小，最小值 8，此时 \(y_{0}=0\)，点 \(P(0,0)\)．}
\fi
\end{ansblock}

\begin{ansblock}[2章-拓-课时16-22]
% ans:2章-拓-课时16-22
\ansitem{23}{D}
\ifansbookdetail
\ansnote{详解}{\(P\) 在抛物线上：\(16=4m\)、\(m=4\)，\(P(4,-4)\)．焦半径（定义）：准线 \(x=-1\)，\(|PF|=x_{P}+1=4+1=5\)．（两点距复核：\(F(1,0)\)，\(|PF|=\sqrt{9+16}=5\) \ding{51}．）选 D．}
\fi
\end{ansblock}

\begin{ansblock}[2章-拓-课时16-06]
% ans:2章-拓-课时16-06
\ansitem{24}{B}
\ifansbookdetail
\ansnote{详解}{由对称性不妨设 \(P(x,y)\)（\(y>0\)）．特殊位：\(x=1\) 时 \(P(1,2)\)，\(PA\perp x\) 轴，\(\tan\angle APB=|AB|/|PA|=2/2=1\)，\(\angle APB=\pi/4\)；\(x=3\) 时 \(P(3,2\sqrt{3})\)，\(PB\perp x\) 轴，\(\tan\angle APB=|AB|/|PB|=2/(2\sqrt{3})=\sqrt{3}/3\)，\(\angle APB=\pi/6\)．一般位 \(x\neq 1\) 且 \(x\neq 3\)：\par\noindent \(\tan\angle APB=(k_{PB}-k_{PA})/(1+k_{PA}\cdot k_{PB})\)，其中 \(k_{PA}=y/(x-1)\)、\(k_{PB}=y/(x-3)\)，代入化简得 \(\tan\angle APB=2y/(x^{2}-4x+3+y^{2})\)；\par\noindent 以 \(y^{2}=4x\) 代 \(x=y^{2}/4\)：\(\tan\angle APB=2y/(y^{4}/16+3)=2/((1/16)y^{3}+3/y)=2/((1/16)y^{3}+1/y+1/y+1/y)\)．\par\noindent 分母四项用均值不等式：\((1/16)y^{3}+1/y+1/y+1/y\geqslant 4\sqrt[4]{(1/16)y^{3}\cdot(1/y)^{3}}=4\cdot(1/16)^{1/4}=2\)，故 \(\tan\angle APB\leqslant 2/2=1\)；等号要求 \((1/16)y^{3}=1/y\) 即 \(y=2\)（\(P(1,2)\)，落在 \(x=1\) 特殊位，不在一般位内），故一般位 \(0<\tan\angle APB<1\)、\(0<\angle APB<\pi/4\)．综上 \(\angle APB\) 最大值 \(\pi/4\)，选 B．}
\fi
\end{ansblock}

\begin{ansblock}[2章-拓-课时16-05]
% ans:2章-拓-课时16-05
\ansitem{25}{D}
\ifansbookdetail
\ansnote{详解}{\(M(x,y)\) 在以 \(B(6,0)\) 为圆心、2 为半径的圆上：\((x-6)^{2}+y^{2}=4\)；\(N(x_{0},y_{0})\) 满足 \(y_{0}^{2}=8x_{0}\)（\(x_{0}\geqslant 0\)）．在圆约束下加权配凑：\par\noindent \((1/4)|MA|^{2}=(1/4)[(x-2)^{2}+y^{2}]=(1/4)[(x-2)^{2}+y^{2}]+(3/4)[(x-6)^{2}+y^{2}-4]=(x-5)^{2}+y^{2}\)（所加项在圆上为 0），即 \((1/2)|MA|=|MC|\)，其中 \(C(5,0)\)．\par\noindent 于是 \(|MN|+(1/2)|MA|=|MN|+|MC|\geqslant|CN|\)．\par\noindent 又 \(|CN|^{2}=(x_{0}-5)^{2}+y_{0}^{2}\)，代入 \(y_{0}^{2}=8x_{0}\) 得 \(|CN|^{2}=(x_{0}-1)^{2}+24\geqslant 24\)，即 \(|CN|\geqslant 2\sqrt{6}\)，当且仅当 \(x_{0}=1\)（\(N(1,\pm 2\sqrt{2})\)）且 \(C\)、\(M\)、\(N\) 三点共线（\(M\) 为线段 \(CN\) 与圆的交点）时取等——取等可行：\(C(5,0)\) 在圆内（\(|CB|=1<2\)）而 \(N\) 在圆外（\(|NB|^{2}=(x_{0}-6)^{2}+8x_{0}=(x_{0}-2)^{2}+32>4\) 恒成立），线段 \(CN\) 必与圆相交．最小值 \(2\sqrt{6}\)，选 D．}
\fi
\end{ansblock}

\begin{ansblock}[2章-拓-课时16-19]
% ans:2章-拓-课时16-19
\ansitem{26}{B}
\ifansbookdetail
\ansnote{详解}{连接 \(PF\)．由抛物线定义，\(P\) 到焦点的距离等于 \(P\) 到准线的距离，即 \(|PF|=|PQ|\)，\(\triangle QPF\) 为等腰三角形，\(P\) 到线段 \(FQ\) 两端点等距，故 \(P\) 在线段 \(FQ\) 的垂直平分线上，选 B．（反例排 A：取 \(y^{2}=4x\)、\(P(1,2)\)，则 \(F(1,0)\)、\(Q(-1,2)\)，\(|OF|=1\) 而 \(|OQ|=\sqrt{5}\)，\(O\) 不在垂直平分线上；C、D 随之不恒成立．）}
\fi
\end{ansblock}

\begin{ansblock}[2章-拓-课时16-11]
% ans:2章-拓-课时16-11
\ansitem{27}{ABC}
\ifansbookdetail
\ansnote{详解}{准线交 \(x\) 轴于 \(P\)，\(|PF|=p\)．分别过 \(A\)、\(B\) 作准线的垂线，垂足 \(E\)、\(M\)．倾斜角 \(60^{\circ}\)：\(AE\parallel x\) 轴 \(\Rightarrow\angle EAF=60^{\circ}\)；由定义 \(|AE|=|AF|\) \(\Rightarrow\triangle AEF\) 为等边三角形 \(\Rightarrow\angle EFP=\angle AEF=60^{\circ}\)、\(\angle FEP=30^{\circ}\)（直角 \(\triangle EPF\) 中），于是 \(|AF|=|EF|=2|PF|=2p=8\)，\(p=4\)，A 对．又 \(|AE|=|EF|=2|PF|\) 且 \(PF\parallel AE\) \(\Rightarrow F\) 为 \(AD\) 中点，\(|DF|=|FA|\)，B 对．\(\angle DAE=60^{\circ}\) 的补算：\(\triangle ADE\) 中 \(\angle AED=90^{\circ}\)、\(\angle ADE=30^{\circ}\) \(\Rightarrow|BD|=2|BM|=2|BF|\)（末步用定义 \(|BM|=|BF|\)），C 对．由 \(|BD|=2|BF|\) 得 \(|BF|=(1/3)|DF|=(1/3)|AF|=8/3\neq 4\)，D 错．故选 ABC．}
\fi
\end{ansblock}

\begin{ansblock}[2章-拓-课时16-12]
% ans:2章-拓-课时16-12
\ansitem{28}{A}
\ifansbookdetail
\ansnote{详解}{圆化标准形：\((x-1)^{2}+y^{2}=1\)，圆心 \(F(1,0)\)、半径 1，恰为抛物线 \(y^{2}=4x\) 的焦点．设 \(A(x_{1},y_{1})\)、\(D(x_{2},y_{2})\)（\(A\)、\(D\) 为直线与抛物线的两交点，\(B\)、\(C\) 为与圆的两交点，自上而下顺次），由焦半径 \(|AF|=x_{1}+1\)、\(|DF|=x_{2}+1\)，得 \(|AB|=|AF|-|BF|=x_{1}+1-1=x_{1}\)，\(|CD|=|DF|-|CF|=x_{2}\)．设 \(l\)：\(x=my+1\) 代入 \(y^{2}=4x\) 得 \(y^{2}-4my-4=0\)，\(y_{1}y_{2}=-4\)，故 \par\noindent \(|AB|\cdot|CD|=x_{1}x_{2}=(y_{1}^{2}/4)(y_{2}^{2}/4)=(y_{1}y_{2})^{2}/16=16/16=1\)，\par\noindent 为定值 1（与 \(m\) 无关，最值之说随之落空）．选 A．}
\fi
\end{ansblock}

\begin{ansblock}[2章-拓-课时16-13]
% ans:2章-拓-课时16-13
\ansitem{29}{B（−4）}
\ifansbookdetail
\ansnote{详解}{联立 \(x^{2}=4(kx+2)\)：\(x^{2}-4kx-8=0\)（\(\Delta=16k^{2}+32>0\) 恒成立），\(x_{1}+x_{2}=4k\)、\(x_{1}x_{2}=-8\)．\par\noindent \(y_{1}y_{2}=(kx_{1}+2)(kx_{2}+2)=k^{2}x_{1}x_{2}+2k(x_{1}+x_{2})+4=-8k^{2}+8k^{2}+4=4\)（即过 \((0,2)\) 的弦两端纵坐标之积为定值 4，\(A\)、\(B\) 分居 \(y\) 轴两侧故 \(x_{1}x_{2}<0\) 与 \(-8\) 相符）．\par\noindent \(\overrightarrow{OA}\cdot\overrightarrow{OB}=x_{1}x_{2}+y_{1}y_{2}=-8+4=-4\)，与 \(k\) 无关．值 \(-4\) 对应选项 B，选 B．}
\fi
\end{ansblock}

\begin{ansblock}[2章-拓-课时16-26]
% ans:2章-拓-课时16-26
\ansitem{30}{A}
\ifansbookdetail
\ansnote{详解}{\(F(1,0)\)．设 \(l\)：\(x=ty+1\)（含斜率不存在），代入 \(y^{2}=4x\)：\(y^{2}-4ty-4=0\)，\(y_{1}+y_{2}=4t\)、\(y_{1}y_{2}=-4\)，\(x_{1}+x_{2}=t(y_{1}+y_{2})+2=4t^{2}+2\)、\(x_{1}x_{2}=(y_{1}y_{2})^{2}/16=1\)．\(\angle AMB\) 为锐角 \(\iff\overrightarrow{MA}\cdot\overrightarrow{MB}>0\)（且 \(A\)、\(M\)、\(B\) 不共线，\(M\) 在 \(x\) 轴负半轴而弦过 \(F\) 恒成立）：\par\noindent \(\overrightarrow{MA}\cdot\overrightarrow{MB}=(x_{1}-m)(x_{2}-m)+y_{1}y_{2}=x_{1}x_{2}-m(x_{1}+x_{2})+m^{2}-4=m^{2}-2m-3-4mt^{2}\)．\par\noindent 按题意，对一切过 \(F\) 的弦（\(\forall t\in\mathrm{R}\)）均存在这样的锐角：\(m<0\) 时 \(-4mt^{2}\geqslant 0\) 且随 \(t^{2}\) 递增，最严处 \(t=0\)，须 \(m^{2}-2m-3>0\) \(\Rightarrow m<-1\) 或 \(m>3\)，配 \(m<0\) 得 \(m\in(-\infty,-1)\)，选 A．}
\fi
\end{ansblock}

\begin{ansblock}[2章-拓-课时16-27]
% ans:2章-拓-课时16-27
\ansitem{31}{B}
\ifansbookdetail
\ansnote{详解}{\(p=4\)，准线 \(x=-2\)．分别过 \(A\)、\(M\)、\(B\) 作准线的垂线，垂足 \(C\)、\(D\)、\(E\)，则 \(MD\) 为梯形（退化情形为中点公式仍合）\(ACEB\) 的中位线，\(|MD|=(|AC|+|BE|)/2=(|AF|+|BF|)/2=|AB|/2\)（定义）．设 \(l\)：\(x=my+2\) 代入 \(y^{2}=8x\)：\(y^{2}-8my-16=0\)，\par\noindent \(|AB|=\sqrt{1+m^{2}}\,|y_{1}-y_{2}|=\sqrt{1+m^{2}}\cdot\sqrt{64m^{2}+64}=8(1+m^{2})\geqslant 8\)（\(m=0\) 即通径取等）．\par\noindent 故 \(|MD|=|AB|/2\geqslant 4\)，最小值 4，选 B．}
\fi
\end{ansblock}

\begin{ansblock}[2章-拓-课时16-28]
% ans:2章-拓-课时16-28
\ansitem{32}{A}
\ifansbookdetail
\ansnote{详解}{由定义 \(d_{1}=|MF|\)；设 \(MN\perp l'\) 于 \(N\)，则 \(d_{2}=|MN|\)，\(d_{1}+d_{2}=|MF|+|MN|\geqslant d(F,l')\)（\(F\) 到 \(l'\) 的垂线段长），当 \(M\) 落在「过 \(F\) 垂直于 \(l'\) 的直线」与抛物线的交点上（该垂线与 \(l'\) 的垂足 \(N\) 在 \(F\) 同侧射影之间，可行）时取最小．\par\noindent \(d(F,l')=|1-0+2|/\sqrt{2}=3\sqrt{2}/2\)．\par\noindent 设 \(l'\) 与 \(x\) 轴交于 \(D(-2,0)\)，\(|FD|=3\)．\par\noindent 在 Rt\(\triangle DNF\)（\(\angle DNF=90^{\circ}\)）中，\(\cos\angle NFD=|FN|/|FD|=(3\sqrt{2}/2)/3=\sqrt{2}/2\)；取最小值时 \(M\) 在线段 \(FN\) 上、射线 \(FO\) 与 \(FD\) 同向，\(\cos\angle MFO=\sqrt{2}/2\)，选 A．（\(M\) 坐标核：\(M(t^{2},2t)\) 且 \(FM\parallel(1,-1)\) \(\Rightarrow t=\sqrt{2}-1\)，\(M(3-2\sqrt{2},2\sqrt{2}-2)\) 在 \(FN\) 段内 \ding{51}．）}
\fi
\end{ansblock}

\begin{ansblock}[2章-拓-课时16-24]
% ans:2章-拓-课时16-24
\ansitem{33}{C}
\ifansbookdetail
\ansnote{详解}{\(|OA|=|OB|\)：\(x_{1}^{2}+y_{1}^{2}=x_{2}^{2}+y_{2}^{2}\)，而 \(y_{i}=x_{i}^{2}/2\)，\((x_{1}^{2}-x_{2}^{2})[1+(x_{1}^{2}+x_{2}^{2})/4]=0\) \(\Rightarrow x_{1}^{2}=x_{2}^{2}\)，\(A\)、\(B\) 异点故 \(x_{2}=-x_{1}\)，两点关于 \(y\) 轴对称．设 \(B(a,a^{2}/2)\)、\(A(-a,a^{2}/2)\)（\(a>0\)），\(S=(1/2)\cdot 2a\cdot(a^{2}/2)=a^{3}/2=12\sqrt{3}\) \(\Rightarrow a^{3}=24\sqrt{3}=(2\sqrt{3})^{3}\) \(\Rightarrow a=2\sqrt{3}\)，\(B(2\sqrt{3},6)\)．记 \(OB\) 与 \(y\) 轴夹角 \(\theta\)：\(\tan\theta=a/(a^{2}/2)=2\sqrt{3}/6=\sqrt{3}/3\) \(\Rightarrow\theta=30^{\circ}\)，\(\angle AOB=2\theta=60^{\circ}\)．选 C．}
\fi
\end{ansblock}

\begin{ansblock}[2章-拓-课时16-15]
% ans:2章-拓-课时16-15
\ansitem{34}{y²=(p/2)x（x≥0），轨迹是顶点在原点、对称轴为 x 轴、开口向右的抛物线（2p̃=p/2，焦点 (p/8,0)）}
\ifansbookdetail
\ansnote{详解}{设抛物线上点 \((x_{0},y_{0})\)，\(y_{0}^{2}=2px_{0}\)（\(x_{0}\geqslant 0\)），垂线段中点 \((x,y)=(x_{0},y_{0}/2)\)．消参：\(x_{0}=x\)、\(y_{0}=2y\) 代入得 \(4y^{2}=2px\)，即 \(y^{2}=(p/2)x\)（\(x\geqslant 0\)）．这是顶点在原点、对称轴为 \(x\) 轴、开口向右的抛物线，其 \(2\tilde{p}=p/2\)（\(\tilde{p}=p/4\)），焦点 \((p/8,0)\)．（原点处垂线段退化为点 \(O\)，中点仍为 \(O\)，轨迹含顶点 \ding{51}．）}
\fi
\end{ansblock}

\begin{ansblock}[2章-拓-课时16-16]
% ans:2章-拓-课时16-16
\ansitem{35}{y²=8x}
\ifansbookdetail
\ansnote{详解}{点 \(A\) 横坐标 \(2>0\) 在 \(y\) 轴右侧，故抛物线开口向右，设 \(y^{2}=2px\)（\(p>0\)），\(F(p/2,0)\)．设 \(A(2,y_{A})\)，\(y_{A}^{2}=4p\)．\(\overrightarrow{FA}=(2-p/2,y_{A})\)、\(\overrightarrow{OA}=(2,y_{A})\)，\(\overrightarrow{FA}\cdot\overrightarrow{OA}=2(2-p/2)+y_{A}^{2}=4-p+4p=4+3p=16\)，解得 \(p=4\)，抛物线为 \(y^{2}=8x\)．（回验：\(A(2,\pm 4)\)，\(\overrightarrow{FA}\cdot\overrightarrow{OA}=2(2-2)+16=16\) \ding{51}．）}
\fi
\end{ansblock}

\begin{ansblock}[2章-拓-课时16-17]
% ans:2章-拓-课时16-17
\ansitem{36}{证明见详解；类似结论：y²=−2px 上 |MF|=p/2−x₀（x₀≤0）；x²=2py 上 |MF|=y₀+p/2；x²=−2py 上 |MF|=p/2−y₀}
\ifansbookdetail
\ansnote{详解}{证明：准线 \(l\)：\(x=-p/2\)．由抛物线定义，\(|MF|\) 等于 \(M\) 到 \(l\) 的距离；\(M\) 横坐标 \(x_{0}\geqslant 0\)，故距离为 \(x_{0}-(-p/2)=x_{0}+p/2\)．类似结论（同法：焦半径＝到准线距离，注意坐标符号范围）：①\(y^{2}=-2px\)（\(x_{0}\leqslant 0\)，准线 \(x=p/2\)）：\(|MF|=p/2-x_{0}\)；②\(x^{2}=2py\)（\(y_{0}\geqslant 0\)，准线 \(y=-p/2\)）：\(|MF|=y_{0}+p/2\)；③\(x^{2}=-2py\)（\(y_{0}\leqslant 0\)，准线 \(y=p/2\)）：\(|MF|=p/2-y_{0}\)．统一记法：焦半径＝\(p/2\)＋（沿开口方向的点到顶点有向距离）．}
\fi
\end{ansblock}

\begin{ansblock}[2章-拓-课时16-07]
% ans:2章-拓-课时16-07
\ansitem{37}{y²=3x 或 y²=−3x}
\ifansbookdetail
\ansnote{详解}{设方程为 \(y^{2}=2px\) 或 \(y^{2}=-2px\)（\(p>0\)），交点 \(A(x_{1},y_{1})\)、\(B(x_{2},y_{2})\)（\(y_{1}>0>y_{2}\)）．两曲线均关于 \(x\) 轴对称，公共弦被 \(x\) 轴垂直平分：\(x_{1}=x_{2}\)、\(y_{2}=-y_{1}\)，弦长 \(y_{1}-y_{2}=2y_{1}=2\sqrt{3}\) \(\Rightarrow y_{1}=\sqrt{3}\)．代入圆：\(x_{1}^{2}=4-3=1\)、\(x_{1}=\pm 1\)．点 \((1,\sqrt{3})\) 在 \(y^{2}=2px\) 上：\(3=2p\)、\(p=3/2\) \(\Rightarrow y^{2}=3x\)；点 \((-1,\sqrt{3})\) 在 \(y^{2}=-2px\) 上：同理 \(p=3/2\) \(\Rightarrow y^{2}=-3x\)．故方程为 \(y^{2}=3x\) 或 \(y^{2}=-3x\)．}
\fi
\end{ansblock}

\begin{ansblock}[2章-拓-课时16-08]
% ans:2章-拓-课时16-08
\ansitem{38}{BCD}
\ifansbookdetail
\ansnote{详解}{圆 \(C\) 圆心 \((0,1)\)、半径 4；抛物线 \(x^{2}=4y\) 的焦点恰为 \(F(0,1)\)、准线 \(y=-1\)，圆心与焦点重合．联立 \(x^{2}=4y\) 与 \(x^{2}+(y-1)^{2}=16\)：\(4y+(y-1)^{2}=16\) \(\Rightarrow y^{2}+2y-15=0\) \(\Rightarrow y=3\)（\(y=-5\) 舍，\(x^{2}=-20\) 无解），交点 \((\pm 2\sqrt{3},3)\)．劣弧 \(AB\) 为弦 \(y=3\) 上方含点 \((0,5)\) 的短弧（圆心到弦距离 \(2<\) 半径 4），其上动点 \(P\) 纵坐标 \(y_{P}\in(3,5)\)，A 项误以横坐标 \(2\sqrt{3}\) 为下界，错．B：\(P\)、\(N\) 同横坐标且 \(P\) 在 \(N\) 上方，\(|PN|=y_{P}-y_{N}\)；由定义 \(|NF|=y_{N}+1\)（\(N\) 到准线距离），故 \(|PN|+|NF|=y_{P}+1\)，恰为 \(P\) 到准线 \(y=-1\) 的距离，对（特位 \(x=0\)：\(P(0,4)\)、\(N(0,0)\)，\(4+1=5=y_{P}+1\) \ding{51}）．C：圆心 \((0,1)\) 到准线 \(y=-1\) 距离为 2，对．D：\(|PF|=\) 半径 4，\(\triangle PFN\) 周长 \(=|PF|+|PN|+|NF|=4+(y_{P}-y_{N})+(y_{N}+1)=y_{P}+5\in(8,10)\)，对．故选 BCD．}
\fi
\end{ansblock}

\begin{ansblock}[2章-拓-课时16-10]
% ans:2章-拓-课时16-10
\ansitem{39}{ABD}
\ifansbookdetail
\ansnote{详解}{A：由定义，\(|PF|\) 等于 \(P\) 到准线 \(l\) 的距离，而 \(PE\perp l\) 于 \(E\)，故 \(|PE|=|PF|\)，对．B：\(PQ\) 平分 \(\angle EPF\) 的外角 \(\Rightarrow\angle FPQ=\angle NPQ\)；又 \(EN\parallel x\) 轴（同垂直于准线）\(\Rightarrow\angle NPQ=\angle PQF\)（内错角），得 \(\angle FPQ=\angle PQF\)，\(\triangle FPQ\) 中 \(|PF|=|QF|\)，对．C：若 \(|PN|=|MF|\)，又 \(|QM|=|QN|\)（\(Q\) 在外角平分线上，到角两边等距）、\(\angle QMF=\angle QNP=90^{\circ}\)，则 \(\triangle FMQ\cong\triangle PNQ\)，得 \(|FQ|=|PQ|\)，结合 B 得 \(\triangle PFQ\) 为正三角形、\(\angle PFQ=\pi/3\)，则 \(P\) 被锁定为定点，与「\(P\) 为异于 \(O\) 的任意一点」矛盾，C 不恒成立．D：四边形 \(KQNE\) 为矩形（\(\angle KEQ=\angle EQN=\angle K=90^{\circ}\)）\(\Rightarrow|EN|=|KQ|\)；\(|PN|=|EN|-|EP|=|KQ|-|PF|=|KF|+|FQ|-|PF|=|KF|\)（用 \(|EP|=|PF|=|FQ|\)），对．故选 ABD．}
\fi
\end{ansblock}

\begin{ansblock}[2章-拓-课时16-14]
% ans:2章-拓-课时16-14
\ansitem{40}{AD}
\ifansbookdetail
\ansnote{详解}{抛物线 \(x^{2}=2y\) 的焦点 \((0,1/2)\)、准线 \(y=-1/2\)．设 \(l\)：\(y=kx+1/2\)，代入 \(x^{2}=2y\) 得 \(x^{2}-2kx-1=0\)，\(x_{1}+x_{2}=2k\)、\(x_{1}x_{2}=-1\)；\(y_{1}+y_{2}=k(x_{1}+x_{2})+1=2k^{2}+1\)、\par\noindent \(y_{1}y_{2}=(kx_{1}+1/2)(kx_{2}+1/2)=k^{2}x_{1}x_{2}+(k/2)(x_{1}+x_{2})+1/4=-k^{2}+k^{2}+1/4=1/4\)．\par\noindent A：过点 \((1,1/2)\)、斜率 \(t\) 的直线 \(y=t(x-1)+1/2\) 代入 \(x^{2}=2y\) 得 \(x^{2}-2tx+2t-1=0\)，即 \((x-1)(x-(2t-1))=0\)，另一交点横标 \(2t-1\)；相切 \(\iff 2t-1=1\iff t=1\)，切线 \(y=x-1/2\) 即 \(2x-2y-1=0\)，A 对．B：\(\overrightarrow{AM}\cdot\overrightarrow{BM}=(-x_{1},-1/2-y_{1})\cdot(-x_{2},-1/2-y_{2})=x_{1}x_{2}+(1/2+y_{1})(1/2+y_{2})=x_{1}x_{2}+1/4+(1/2)(y_{1}+y_{2})+y_{1}y_{2}=-1+1/4+(1/2)(2k^{2}+1)+1/4=k^{2}\)，仅 \(k=0\)（通径位）时为 0，不恒成立，B 错．C：取 \(l\parallel x\) 轴（\(k=0\)，即通径），\(A(1,1/2)\)、\(B(-1,1/2)\)，\(|OA|=|OB|=\sqrt{1+1/4}=\sqrt{5}/2\)，\(|OA|+|OB|=\sqrt{5}\)，严格不等式不成立，C 错．D：\(BN\perp\) 准线 \(\Rightarrow N(x_{2},-1/2)\)；直线 \(OA\)：\(y=(y_{1}/x_{1})x=(x_{1}/2)x\)（用 \(y_{1}=x_{1}^{2}/2\)），代入 \(x=x_{2}\) 得纵标 \((x_{1}x_{2})/2=-1/2\)，即直线 \(OA\) 恰过 \(N(x_{2},-1/2)\)，\(A\)、\(O\)、\(N\) 共线，D 对．故选 AD．}
\fi
\end{ansblock}

\begin{ansblock}[2章-拓-课时16-20]
% ans:2章-拓-课时16-20
\ansitem{41}{2√3}
\ifansbookdetail
\ansnote{详解}{椭圆左顶点 \((-4,0)\)．设 \(P(-y^{2}/4,y)\)，则 \par\noindent \(d^{2}=(-y^{2}/4+4)^{2}+y^{2}=(1/16)y^{4}-y^{2}+16=(1/16)(y^{2}-8)^{2}+12\geqslant 12\)，\par\noindent 当 \(y^{2}=8\)（\(y=\pm 2\sqrt{2}\)，\(x=-2\)，点 \(P(-2,\pm 2\sqrt{2})\) 在抛物线上 \ding{51}）时取等，\(d\) 的最小值为 \(\sqrt{12}=2\sqrt{3}\)．}
\fi
\end{ansblock}

\begin{ansblock}[2章-拓-课时16-23]
% ans:2章-拓-课时16-23
\ansitem{42}{y²=12x}
\ifansbookdetail
\ansnote{详解}{设圆 \(M\) 半径为 \(R\)．\(\odot C\) 圆心 \(C(3,0)\)、半径 1，外切 \(\Rightarrow|MC|=R+1\)；\(\odot M\) 与直线 \(x=-2\) 相切 \(\Rightarrow\) 圆心 \(M\) 到 \(x=-2\) 的距离为 \(R\)，从而 \(M\) 到直线 \(x=-3\) 的距离为 \(R+1\)．于是动点 \(M\) 到定点 \(C(3,0)\) 的距离等于它到定直线 \(x=-3\) 的距离，轨迹是以 \(C\) 为焦点、\(x=-3\) 为准线的抛物线（定义）：\(p/2=3\)、\(p=6\)，标准方程 \(y^{2}=12x\)．（核验 \(x\geqslant 0\)：\(M\) 到 \(x=-3\) 距离 \(=\) 横坐标 \(+3=R+1>0\) 且 \(|MC|=R+1\geqslant 1>0\) 自洽．）}
\fi
\end{ansblock}

\begin{ansblock}[2章-拓-课时16-25]
% ans:2章-拓-课时16-25
\ansitem{43}{7√3/12}
\ifansbookdetail
\ansnote{详解}{等腰梯形两底 \(AB\parallel CD\)，而抛物线上平行于 \(x\) 轴方向的弦（两交点同纵坐标）不存在（\(y^{2}=2px\) 单值），故两底必为垂直于 \(x\) 轴的弦：\(A\)、\(B\) 关于 \(x\) 轴对称，\(C\)、\(D\) 关于 \(x\) 轴对称．\(|AB|=2\Rightarrow y_{A}=1\)（取上半），\(|CD|=4\Rightarrow y_{D}=2\)．设 \(A(m,1)\)；腰与底夹角 \(\angle ADC=60^{\circ}\)，底铅直，故 \(DA\) 与铅直线夹 \(60^{\circ}\)，纵差 \(2-1=1\) \(\Rightarrow\) 横差 \(=\tan 60^{\circ}=\sqrt{3}\)，\(D(m+\sqrt{3},2)\)．两点在抛物线上：\(1=2pm\)、\(4=2p(m+\sqrt{3})\)，两式相减 \(3=2\sqrt{3}p\) \(\Rightarrow p=\sqrt{3}/2\)、\(m=1/(2p)=\sqrt{3}/3\)．由焦半径（例 36 结论）\(|AF|=m+p/2=\sqrt{3}/3+\sqrt{3}/4=7\sqrt{3}/12\)．}
\fi
\end{ansblock}

\begin{ansblock}[2章-拓-课时16-29]
% ans:2章-拓-课时16-29
\ansitem{44}{(1) y²=4x；(2) 直线 AB 恒过定点 (−1,0)}
\ifansbookdetail
\ansnote{详解}{(1) 设圆心 \(M(x,y)\)，圆过 \((2,0)\) \(\Rightarrow\) 半径\({}^{2}=(x-2)^{2}+y^{2}\)；被 \(y\) 轴截弦长 4 \(\Rightarrow\) 半径\({}^{2}=x^{2}+2^{2}\)（半弦 2、弦心距 \(|x|\)）．两式相等：\((x-2)^{2}+y^{2}=x^{2}+4\) \(\Rightarrow y^{2}=4x\)．\par\noindent (2) 顶点在轨迹 \(y^{2}=4x\) 上．设 \(B(x_{1},y_{1})\)、\(C(x_{2},y_{2})\)，\(BC\) 过 \(F(1,0)\) 且不垂直于坐标轴，设 \(BC\)：\(x=ty+1\)：\(y^{2}-4ty-4=0\) \(\Rightarrow y_{1}y_{2}=-4\)．由 \(A\)、\(C\) 关于 \(x\) 轴对称，\(A(x_{2},-y_{2})\)．设 \(AB\)：\(y=kx+m\)（\(k\neq 0\)）代入 \(ky^{2}-4y+4m=0\)，其两根为 \(y_{1}\)、\(-y_{2}\)：\(y_{1}\cdot(-y_{2})=4m/k\)，即 \(4=4m/k\)、\(m=k\)．故 \(AB\)：\(y=k(x+1)\)，对一切实有 \(k\) 恒过定点 \((-1,0)\)．}
\fi
\end{ansblock}

\begin{ansblock}[2章-拓-课时16-30]
% ans:2章-拓-课时16-30
\ansitem{45}{(1) x²=4y；(2) 点 P 恒在定直线 y=−1（即 C 的准线）上}
\ifansbookdetail
\ansnote{详解}{(1) \(A\) 横坐标 \(2\sqrt{2}\)、在圆上：\(y_{A}=(12-8)/2=2\)，\(A(2\sqrt{2},2)\)．代入 \(x^{2}=2py\)：\(8=4p\)、\(p=2\)，抛物线为 \(x^{2}=4y\)．\par\noindent (2) 先给切线式并验证：点 \(M(x_{1},x_{1}^{2}/4)\) 处，直线 \(y=(x_{1}/2)x-x_{1}^{2}/4\) 与 \(x^{2}=4y\) 联立得 \(x^{2}-2x_{1}x+x_{1}^{2}=(x-x_{1})^{2}=0\)，唯一公共点（重根 \(x_{1}\)），即为过 \(M\) 的切线 \(l_{1}\)；同理 \(N(x_{2},x_{2}^{2}/4)\) 处切线 \(l_{2}\)：\(y=(x_{2}/2)x-x_{2}^{2}/4\)．两切线联立：\((x_{1}-x_{2})x/2=(x_{1}^{2}-x_{2}^{2})/4\) \(\Rightarrow x_{0}=(x_{1}+x_{2})/2\)，\(y_{0}=(x_{1}/2)x_{0}-x_{1}^{2}/4=x_{1}x_{2}/4\)．设 \(l\)：\(y=kx+1\)（过 \(F(0,1)\)；\(l\) 为 \(x=0\) 时与抛物线仅一交点，不合「不同两点」），代入 \(x^{2}=4y\)：\(x^{2}-4kx-4=0\) \(\Rightarrow x_{1}x_{2}=-4\)，故 \(y_{0}=-1\) 恒成立，点 \(P\) 恒在定直线 \(y=-1\)（准线）上．}
\fi
\end{ansblock}

\begin{ansblock}[2章-导-课时16-G1]
% ans:2章-导-课时16-G1
\ansitem{46}{A}
\ifansbookdetail
\ansnote{详解}{\(y=2x^{2}\) 即 \(x^{2}=(1/2)y\)，\(2p=1/2\)、\(p=1/4\)，开口向上，焦点 \((0,p/2)=(0,1/8)\)．故选 A．（验算：准线 \(y=-1/8\)，顶点 \((0,0)\) 到焦点、到准线等距 \(1/8\) \ding{51}．）}
\fi
\end{ansblock}

\begin{ansblock}[2章-导-课时16-G2]
% ans:2章-导-课时16-G2
\ansitem{47}{D}
\ifansbookdetail
\ansnote{详解}{\(x^{2}=4y\) 的焦点在 \(y\) 轴上，图象关于 \(y\) 轴对称，对称轴为直线 \(x=0\)．故选 D．（验算：点 \((2,1)\) 在曲线上（\(4=4\times 1\) \ding{51}），其关于 \(y\) 轴的对称点 \((-2,1)\) 也在曲线上 \ding{51}；而关于 \(x\) 轴的对称点 \((2,-1)\) 不在曲线上（\(4\neq -1\)）——对称轴确为 \(y\) 轴 \ding{51}．）}
\fi
\end{ansblock}

\begin{ansblock}[2章-导-课时16-G3]
% ans:2章-导-课时16-G3
\ansitem{48}{D}
\ifansbookdetail
\ansnote{详解}{对圆上一点 \(Q\)，\(|PQ|\geqslant|PM|-|MQ|=|PM|-1\)（\(M\) 为圆心），且当 \(P\)、\(Q\)、\(M\) 共线、\(Q\) 在线段 \(PM\) 上时取等．设 \(P(x_{0},y_{0})\) 在 \(y^{2}=x\) 上，则 \(x_{0}=y_{0}^{2}\)，圆心 \(M(3,0)\)：\par\noindent \(|PM|^{2}=(y_{0}^{2}-3)^{2}+y_{0}^{2}\)．令 \(t=y_{0}^{2}\)（\(t\geqslant 0\)）：\(|PM|^{2}=(t-3)^{2}+t=t^{2}-5t+9=(t-5/2)^{2}+11/4\geqslant 11/4\)，\par\noindent 故 \(|PM|_{\min}=\sqrt{11}/2\)（\(t=5/2\) 时取得，此时 \(y_{0}=\pm\sqrt{10}/2\)，\(P(5/2,\pm\sqrt{10}/2)\) 在抛物线上 \ding{51}），从而 \(|PQ|_{\min}=\sqrt{11}/2-1\)，选 D．（验算：\(t=5/2\)：\(|PM|^{2}=(5/2-3)^{2}+5/2=1/4+5/2=11/4\) \ding{51}；\(|PM|_{\min}=\sqrt{11}/2>1\)，圆内不含抛物线点，\(|PQ|=|PM|-1\) 恒正、等号可取 \ding{51}．）}
\fi
\end{ansblock}

\begin{ansblock}[2章-导-课时16-G4]
% ans:2章-导-课时16-G4
\ansitem{49}{0<a≤1}
\ifansbookdetail
\ansnote{详解}{设 \(P(x,y)\) 为抛物线上任意一点，则 \(y=x^{2}/2\)（\(y\geqslant 0\)）．\par\noindent \(|AP|^{2}=x^{2}+(y-a)^{2}=2y+(y-a)^{2}=y^{2}+2(1-a)y+a^{2}\)．\par\noindent 这是关于 \(y\) 的二次函数，对称轴 \(y=a-1\)、开口向上．「最近点恰为顶点 \((0,0)\)」\(\iff|AP|^{2}\) 在 \(y=0\) 处取最小值 \(\iff\) 对称轴不在开区间 \((0,+\infty)\) 内 \(\iff a-1\leqslant 0\)，即 \(a\leqslant 1\)．又 \(a>0\)，故 \(0<a\leqslant 1\)．（验算：\(a=1\)：\(|AP|^{2}=y^{2}+1\) 在 \(y=0\) 处最小 \ding{51}（顶点最近，边界值取得到）；\(a=1.1\)：对称轴 \(y=0.1>0\)，最近点 \(y=0.1\neq 0\) 非顶点 \ding{55}——故右端取闭 \ding{51}．）}
\fi
\end{ansblock}

\begin{ansblock}[2章-导-课时16-G5]
% ans:2章-导-课时16-G5
\ansitem{50}{3√2/2}
\ifansbookdetail
\ansnote{详解}{设 \(P(t^{2}/4,t)\) 为抛物线上任意一点，则 \par\noindent \(d=|t^{2}/4-t+4|/\sqrt{1^{2}+1^{2}}=|t^{2}/4-t+4|/\sqrt{2}\)．\par\noindent 配方：\(t^{2}/4-t+4=(1/4)(t-2)^{2}+3>0\) 恒成立，故 \(d=(1/4)(t-2)^{2}/\sqrt{2}+3/\sqrt{2}\geqslant 3/\sqrt{2}=3\sqrt{2}/2\)，当 \(t=2\) 即 \(P(1,2)\) 时取得．（验算：\(P(1,2)\)：\(d=|1-2+4|/\sqrt{2}=3/\sqrt{2}=3\sqrt{2}/2\) \ding{51}；联立 \(x=y^{2}/4\) 与直线得 \(y^{2}-4y+16=0\)，\(\Delta=16-64<0\)，直线与抛物线无公共点，最小值为正、在切点型位置取得 \ding{51}．）}
\fi
\end{ansblock}

% ---- 尾块（\tailfill[30mm] 定高参——钉档 §三.3：无参在平衡栏呈「内容尾随框」，贴栏底须定高参；实测无参致末页平衡 Overfull\vbox 12.99pt，定高 30mm 三零） ----
\tailfill[30mm]

\end{multicols}
\end{document}
